\documentclass[12pt]{article}
\usepackage[utf8]{inputenc}
\usepackage[T1]{fontenc}
\usepackage{hyperref}
\title{Tomedes Language Resources for K-12 Curriculum Development}
\author{Compiled by Claude Sonnet 5 (claude-sonnet-5) for Andr\'as Kornai}
\date{}
\begin{document}
\maketitle

\section{Introduction}
The ISO 639-3 code toe (Tomedes) was retired by Change Request 2015-033, effective 2016-01-15, as a non-existent language. The IANA subtag was deprecated on 2016-05-30. See \url{https://www.iana.org/assignments/lang-subtags-templates/toe.txt}. Glottolog still lists the entry tome1239: \url{https://glottolog.org/resource/languoid/id/tome1239.bigmap.html}. There is no evidence of a speaker community.

\section{Primary Educational Resources}
None exist.

\section{Linguistic and Reference Materials}
None found.

\section{Religious and Community Resources}
None found.

\section{Digital and Audio Resources}
None found.

\section{Curriculum Development Context}
Because the code names a non-existent language, no K-12 curriculum can be built for it. I suggest removing it from the survey.
\end{document}